\documentclass[11pt]{article}
\usepackage{amsmath}
\usepackage{amssymb}
\usepackage{mathpazo}
\usepackage[dvipsnames]{xcolor}
\usepackage[colorlinks, linkcolor=Blue]{hyperref}
\usepackage{graphicx}
\usepackage[margin=1in]{geometry}
\title{Homework 10}
\author{Your name here}
\date{\today{}}
\newcommand{\w}{\mathbf{w}}
\newcommand{\what}{\widehat{\w}}

\begin{document}
\maketitle

\section{Predictive variance (error bars)}

Include the plots generated by \texttt{hw.py} showing the predictive variance
for the polynomial regression models with orders 1, 3, 5, and 9.  Add a caption to
the figure that qualitatively describes what these plots show, contrasts them
with each other, and explains the effect of removing the points $2.5\leq x \leq
4$ from the initial data.

\begin{figure}[htbp]
    \centering
    %\includegraphics[width=0.4\textwidth]{images/predictive_variance_order_1.pdf}
    %\includegraphics[width=0.4\textwidth]{images/predictive_variance_order_3.pdf}
    %\includegraphics[width=0.4\textwidth]{images/predictive_variance_order_5.pdf}
    %\includegraphics[width=0.4\textwidth]{images/predictive_variance_order_9.pdf}
%%% Answer START %%%
    \caption{\textcolor{RoyalBlue}{Plots of fitting polynomial functions of orders 1, 3, 5, and 9 to
        the data, and plotting error bars for a set of new input-prediction
        pairs.
        For all of the plots, the error bars get larger at the edges of the
        data.
        The order-1 polynomial is unable to capture the deterministic trend in
        the data, and assumes that most of the variability of the data to be
        noise. The error bars are fairly large, and are not very much affected
        by the gap in the data.
        Note the error bars get significantly
        larger in the interval where the data is missing for orders 3, 5, and 9.
        Visually, it seems that the 3rd order polynomial has the smallest error
        bars.}}
    \label{fig:error_bars}
%%% Answer END %%%
\end{figure}

\newpage

\section{Sampling functions}

Include the figures showing the functions corresponding to different values of
$\mathbf{w}$ sampled from $\mathcal{N}(\widehat{\mathbf{w}}, \textsf{cov}\{\widehat{\mathbf{w}}\})$.
Add a caption to the figure that qualitatively describes what these
plots show, contrasts them with each other, and explains the effect of removing
the points $2.5\leq x \leq 4$ from the initial data.

\begin{figure}[htbp]
    \centering
    %\includegraphics[width=0.4\textwidth]{images/sampled_functions_order_1.pdf}
    %\includegraphics[width=0.4\textwidth]{images/sampled_functions_order_3.pdf}\\
    %\includegraphics[width=0.4\textwidth]{images/sampled_functions_order_5.pdf}
    %\includegraphics[width=0.4\textwidth]{images/sampled_functions_order_9.pdf}
%%% Answer START %%%
    \caption{\textcolor{RoyalBlue}{
        Plots of fitting polynomial functions of orders 1, 3, 5, and 9 to
        the data, where the parameters $\w{}$ of each function were
        sampled from $\mathcal{N}(\what{}, \textsf{cov}\{\what\})$ of the best-fit model to the data. As with
        the error bars, we see that the sampled functions vary more in the
        interval where the data was missing, but also there is a definite
        increase in variability for higher-order functions. Again, the model
        of order 3 appears to have the least variance.
    }}
    \label{fig:sampling_functions}
%%% Answer END %%%
\end{figure}

\newpage

\section{Sampling datasets}

Include the figures showing the functions corresponding to different datasets
sampled from polynomial models of orders 1, 3, 5, and 9 with noise. 
Add a descriptive caption to the figure below, and describe what happens to
the variability in possible functions as we change the polynomial order.

\begin{figure}[htbp]
    \centering
    %\includegraphics[width=0.4\textwidth]{images/sampled_datasets_order_1.pdf}
    %\includegraphics[width=0.4\textwidth]{images/sampled_datasets_order_3.pdf}\\
    %\includegraphics[width=0.4\textwidth]{images/sampled_datasets_order_5.pdf}
    %\includegraphics[width=0.4\textwidth]{images/sampled_datasets_order_9.pdf}
%%% Answer START %%%
    \caption{\textcolor{RoyalBlue}{
        Plots of fitting polynomial functions of varying order (1, 3, 5, and 9)
        to 20 independent samples of data. The mean of the true generating
        function is plotted as a darker black line, the other functions are in
        lighter gray lines.
        The plots of model orders 3, 5, and 9 show that as model order
        increases, the best fit line becomes more sensitive to the
        specific data, a sign of over-fitting.
        We can also see that in model orders 5 and 9, the variance increases at
        the boundaries of the data set ranges.
        }}
    \label{fig:sampling_functions}
%%% Answer END %%%
\end{figure}

\newpage

\section{Workload}

How much time did you spend on this homework assignment outside of class?

\section*{Acknowledgments}

Cite all the people you've worked with on this homework, as well as any other
resources you used apart from the FCML textbook.

\end{document}
